
On a alors, d'après IV (2.12.4) :
\[
f_*\left(\sum_{i=0}^{n-1}\xi^{-i}f^*(x_i)\right)
=\sum_{i=0}^{n-1}x_i f_*(\xi^{-i})=x_0,
\]
la dernière égalité résultant de ce que pour $-n+1\le i<0$, on a $f_*(\xi^i)=0$ (VI 2.8). Supposons que $\varphi((x_i))=0$. Alors ce qui précède montre que $x_0=0$; appliquant le même raisonnement aux éléments $\xi^i\varphi((x_i))$ pour $i=1,\ldots,n-1$, on en déduit par récurrence que les $x_i$ sont tous nuls, d'où l'injectivité.

3.3. Soient toujours $S$ un schéma noethérien, et $\underline{E}$ un $\underline{O}_S$-Module localement libre de rang $n$. On note $V=\mathbb V(\underline{E})$ le fibré vectoriel (contravariant en $E$!) associé à $\underline{E}$ (SGA 3 I 4.6), considéré comme schéma en groupes sur $S$. Soit $P$ un torseur sous $V$; on se propose de calculer $K_\bullet(P)$.

Rappelons comment on peut plonger $P$ dans un fibré projectif sur $S$. Les torseurs sous $V$ correspondent aux extensions $\underline{F}$ de $\underline{E}$ par $\underline{O}_S$. En effet, si on a une extension :
\[
0\longrightarrow \underline{O}_S\xrightarrow{s}\underline{F}\xrightarrow{u}\underline{E}\longrightarrow 0,
\tag{3.3.1}
\]
on en déduit un torseur sous $V$ en considérant l'homomorphisme $\mathbb V(\underline{F})\to \mathbb V(\underline{O}_S)$ défini par $s$, la section de $\mathbb V(\underline{O}_S)=\Spec(\underline{O}_S[T])$ définie par $T=1$, et en prenant $P=\mathbb V(\underline{F})\times_{\mathbb V(\underline{O}_S)}S$, les opérations de $V$ sur $P$ étant induites par celles de $V$ sur $\mathbb V(\underline{F})$.

En sens inverse, soit $P$ un torseur sous $V$. On note $\underline P$ le faisceau des germes d'homomorphismes de $S$ dans $P$; c'est un faisceau principal homogène sous le faisceau des germes d'homomorphismes de $S$ dans $V$, qui s'identifie au dual $\check{\underline{E}}$ de $\underline{E}$. On en déduit alors de façon classique une extension de $\underline{O}_S$ par $\check{\underline{E}}$, d'où en passant aux duaux la suite (3.3.1).

